\documentclass[11pt,a4paper]{scrartcl}

\usepackage{amssymb}
\usepackage{amsmath}
\usepackage{tikz}
\usetikzlibrary{angles,quotes,calc,decorations.markings,intersections}

\usepackage[utf8]{inputenc}
\usepackage[T1]{fontenc}
\newcommand{\ats}[1]{\par\noindent\textit{Atsakymas:} #1}
\usepackage{cancel}
\usepackage{tikz}
\usepackage{lmodern} 
\usepackage{amsmath,amssymb,amsthm}
\usepackage{graphicx}
\usepackage[dvipsnames,svgnames]{xcolor}
\usepackage{hyperref}
\usepackage{mathrsfs}
\usepackage[shortlabels]{enumitem}
\usepackage[obeyFinal,textsize=scriptsize,shadow,loadshadowlibrary]{todonotes}
\usepackage{mathtools}
\usepackage{microtype}

%%% Paragraph layout
\setlength{\parindent}{0pt}
\setlength{\parskip}{0.6\baselineskip}

%%% Colored sections
\renewcommand*{\sectionformat}%
  {\color{purple}\S\thesection\enskip}
\renewcommand*{\subsectionformat}%
  {\color{purple}\S\thesubsection\enskip}
\renewcommand*{\subsubsectionformat}%
  {\color{purple}\S\thesubsubsection\enskip}
\addtokomafont{paragraph}{\color{orange!35!black}\P\ }
\KOMAoptions{numbers=noenddot}

%%% Title
\addtokomafont{subtitle}{\Large}
\setkomafont{author}{\Large\scshape}
\setkomafont{date}{\Large\normalsize}
% Neigiamas vertikalus tarpas, kuris pakelia dokumento antraštę į viršų. 
% Galite keisti -2.5cm į -3cm (aukščiau) arba -1.5cm (žemiau).
\titlehead{\vspace*{-7.5cm}} 

% PRIDĖTA: Automatiškai perrašome datą į mokyklos pavadinimą
\AtBeginDocument{%
  \date{\normalsize Vilniaus licėjus}
}

%%% Page setup
\usepackage[headsepline]{scrlayer-scrpage}
\renewcommand{\headfont}{}
\addtolength{\textheight}{3.14cm}
\setlength{\footskip}{0.5in}
\setlength{\headsep}{10pt}
% Puslapio antraštė turi rodyti tik konkrečius dokumento duomenis.
% Nustatome juos aiškiai, kad neatsirastų "\@author" / "\@title" arba senų reikšmių.
\makeatletter
\AtBeginDocument{%
  \ihead{\footnotesize\textbf{Andrius Berniukevičius}}%
  \ohead{\footnotesize\textbf{\@title}}%
}
\makeatother
\automark{section}
\chead{}
\cfoot{\pagemark}

%%% Macros
\providecommand{\ol}{\overline}
\providecommand{\ul}{\underline}
\providecommand{\wt}{\widetilde}
\providecommand{\wh}{\widehat}
\providecommand{\eps}{\varepsilon}
\providecommand{\half}{\frac{1}{2}}
\providecommand{\inv}{^{-1}}
\newcommand{\dang}{\measuredangle} %% Directed angle
\newcommand{\vocab}[1]{\textbf{\color{blue}\sffamily #1}}
\providecommand{\CC}{\mathbb C}
\providecommand{\FF}{\mathbb F}
\providecommand{\NN}{\mathbb N}
\providecommand{\QQ}{\mathbb Q}
\providecommand{\RR}{\mathbb R}
\providecommand{\ZZ}{\mathbb Z}
\providecommand{\ts}{\textsuperscript}
\providecommand{\dg}{^\circ}
\providecommand{\ii}{\item}
\providecommand{\defeq}{\coloneqq}
\DeclareMathOperator*{\lcm}{lcm}
\DeclareMathOperator*{\argmin}{arg min}
\DeclareMathOperator*{\argmax}{arg max}
\providecommand{\hrulebar}{\par
\hspace{\fill}\rule{0.95\linewidth}{.7pt}\hspace{\fill}
\par\nointerlineskip \vspace{\baselineskip}}

%%% Optional colored theorems
\usepackage{thmtools}
\usepackage[framemethod=TikZ]{mdframed}
\mdfdefinestyle{mdbluebox}{%
  roundcorner=10pt,
  linewidth=1pt,
  skipabove=12pt,
  innerbottommargin=9pt,
  skipbelow=2pt,
  linecolor=blue,
  nobreak=true,
  backgroundcolor=TealBlue!5,
}
\declaretheoremstyle[
  headfont=\sffamily\bfseries\color{MidnightBlue},
  mdframed={style=mdbluebox},
  headpunct={\\[3pt]},
  postheadspace={0pt}
]{thmbluebox}
\mdfdefinestyle{mdredbox}{%
  linewidth=0.5pt,
  skipabove=12pt,
  frametitleaboveskip=5pt,
  frametitlebelowskip=0pt,
  skipbelow=2pt,
  frametitlefont=\bfseries,
  innertopmargin=4pt,
  innerbottommargin=8pt,
  nobreak=true,
  backgroundcolor=Salmon!5,
  linecolor=RawSienna,
}
\declaretheoremstyle[
  headfont=\bfseries\color{RawSienna},
  mdframed={style=mdredbox},
  headpunct={\\[3pt]},
  postheadspace={0pt},
]{thmredbox}
\mdfdefinestyle{mdgreenbox}{%
  skipabove=8pt,
  linewidth=2pt,
  rightline=false,
  leftline=true,
  topline=false,
  bottomline=false,
  linecolor=ForestGreen,
  backgroundcolor=ForestGreen!5,
}
\declaretheoremstyle[
  headfont=\bfseries\sffamily\color{ForestGreen!70!black},
  bodyfont=\normalfont,
  spaceabove=2pt,
  spacebelow=1pt,
  mdframed={style=mdgreenbox},
  headpunct={ --- },
]{thmgreenbox}
\mdfdefinestyle{mdblackbox}{%
  skipabove=8pt,
  linewidth=3pt,
  rightline=false,
  leftline=true,
  topline=false,
  bottomline=false,
  linecolor=black,
  backgroundcolor=RedViolet!5!gray!5,
}
\declaretheoremstyle[
  headfont=\bfseries,
  bodyfont=\normalfont\small,
  spaceabove=0pt,
  spacebelow=0pt,
  mdframed={style=mdblackbox}
]{thmblackbox}

%% Aplinkos su rėmeliais (thmtools)
\declaretheorem[style=thmbluebox,name=Teorema]{theorem}
\declaretheorem[style=thmbluebox,name=Lema]{lemma}
\declaretheorem[style=thmbluebox,name=Savybė]{property}
\declaretheorem[style=thmbluebox,name=Teiginys]{proposition}
\declaretheorem[style=thmbluebox,name=Išvada]{corollary}
\declaretheorem[style=thmbluebox,name=Teorema,numbered=no]{theorem*}
\declaretheorem[style=thmbluebox,name=Lema,numbered=no]{lemma*}
\declaretheorem[style=thmbluebox,name=Teiginys,numbered=no]{proposition*}
\declaretheorem[style=thmbluebox,name=Išvada,numbered=no]{corollary*}
\declaretheorem[style=thmgreenbox,name=Algoritmas]{algorithm}
\declaretheorem[style=thmgreenbox,name=Algoritmas,numbered=no]{algorithm*}
\declaretheorem[style=thmgreenbox,name=Teiginys]{claim}
\declaretheorem[style=thmgreenbox,name=Teiginys,numbered=no]{claim*}
\declaretheorem[style=thmredbox,name=Pavyzdys]{example}
\declaretheorem[style=thmredbox,name=Pavyzdys,numbered=no]{example*}
\declaretheorem[style=thmblackbox,name=Pastaba]{remark}
\declaretheorem[style=thmblackbox,name=Pastaba,numbered=no]{remark*}

%% Standartinės amsthm aplinkos (Apibrėžimai ir užduotys)
\theoremstyle{definition}
\ifcsname conjecture\endcsname\else\newtheorem{conjecture}{Hipotezė}\fi
\ifcsname definition\endcsname\else\newtheorem{definition}{Apibrėžimas}\fi
\ifcsname fact\endcsname\else\newtheorem{fact}{Faktas}\fi
\ifcsname answer\endcsname\else\newtheorem{answer}{Atsakymas}\fi
\ifcsname ques\endcsname\else\newtheorem{ques}{Klausimas}\fi
\ifcsname exercise\endcsname\else\newtheorem{exercise}{Užduotis}\fi
\ifcsname problem\endcsname\else\newtheorem{problem}{Uždavinys}[section]\fi
\ifcsname conjecture*\endcsname\else\newtheorem*{conjecture*}{Hipotezė}\fi
\ifcsname definition*\endcsname\else\newtheorem*{definition*}{Apibrėžimas}\fi
\ifcsname fact*\endcsname\else\newtheorem*{fact*}{Faktas}\fi
\ifcsname answer*\endcsname\else\newtheorem*{answer*}{Atsakymas}\fi
\ifcsname ques*\endcsname\else\newtheorem*{ques*}{Klausimas}\fi
\ifcsname exercise*\endcsname\else\newtheorem*{exercise*}{Užduotis}\fi
\ifcsname problem*\endcsname\else\newtheorem*{problem*}{Uždavinys}\fi

\newcounter{answerssection}
\newcommand{\answerssection}{%
  \setcounter{answerssection}{\value{section}}%
  \section{Atsakymai}%
  \setcounter{problem}{0}%
  \renewcommand{\theproblem}{\arabic{answerssection}.\arabic{problem}}%
}

%% GALUTINIS SPRENDIMAS PROOF APLINKAI
%% --- PRADŽIA: Įrodymo aplinkos sutvarkymas ---
\makeatletter

% 1. Užtikriname, kad žodis būtų "Įrodymas", net jei "babel" paketas bando jį perrašyti
\AtBeginDocument{%
  \renewcommand{\proofname}{Įrodymas}%
  \@ifpackageloaded{babel}{%
    \addto\captionslithuanian{\renewcommand{\proofname}{Įrodymas}}%
  }{}%
}

% 2. Priverstinai pašaliname scrartcl klasės bold šriftą ir paliekame TIK italic
% 3. Sujungiame su tuščiu kvadratėliu dešiniajame kampe.
\renewcommand{\qedsymbol}{\ensuremath{\Box}}
\renewenvironment{proof}[1][\proofname]{\par
  \pushQED{\qed}%
  \normalfont \topsep6\p@\@plus6\p@\relax
  \trivlist
  \item[\hskip\labelsep
        \normalfont\itshape % Priverčia naudoti tik pasvirą šriftą, kaip nuotraukoje
    #1\@addpunct{.}]\ignorespaces
}{%
  \popQED\endtrivlist\@endpefalse
}
\newenvironment{solution}{\par
  \normalfont \topsep6\p@\@plus6\p@\relax
  \trivlist
  \item[\hskip\labelsep
        \normalfont\itshape
    \textit{Sprendimas}\@addpunct{.}]\ignorespaces
}{%
  \endtrivlist\@endpefalse
}
\newcommand{\ans}[1]{\par\noindent\textit{Atsakymas:} #1}
\makeatother


\title{Geometrija: kampų medžioklės pradžiamokslis}
\author{Andrius Berniukevičius}

\newenvironment{solution}
    {\begin{list}{}{\setlength{\leftmargin}{10pt}\setlength{\rightmargin}{10pt}}\item[]}
    {\end{list}}

\begin{document}

\maketitle

\section{Naujos žaidimo taisyklės}

Mokykloje geometrijos uždaviniai dažniausiai reikalauja ką nors apskaičiuoti: „duota atkarpa 5, atkarpa 4, raskite plotą“. Olimpiadose tokio tipo geometrijos beveik nėra. Čia svarbiausia – \textbf{struktūra ir įrodymai}.

Vienas svarbiausių mūsų įrankių bus \textbf{Kampų medžioklė} (\textit{Angle Chasing}). 
Kokia šio metodo esmė? \textbf{Mes nenaudojame skaičių!} Jei matome nežinomą kampą, pavadiname jį raide $\alpha$ (alfa), $\beta$ (beta) arba $\gamma$ (gama) ir pradedame „vaikščioti“ po brėžinį, perkeldami šias raides pagal griežtas taisykles, kol atrandame ryšius tarp iš pirmo žvilgsnio nesusijusių figūrų.

% ==========================================
% BAZINIAI GINKLAI
% ==========================================
\section{Medžiotojo arsenalas: 3 pagrindiniai ginklai}

Norint sekti kampus, nereikia dešimčių formulių. Visi olimpiadiniai uždaviniai be apskritimų remiasi vos trimis mechanizmais.

\subsection*{1. Klonavimo aparatas: Lygiašonis trikampis}
Lygiašonis trikampis yra pats svarbiausias kampų perkėlimo įrankis. Jei matote (arba įrodote), kad dvi atkarpos lygios, jūs akimirksniu \textbf{klonuojate} kampą $\alpha$ į kitą trikampio pusę.
\begin{center}
\begin{tikzpicture}[scale=1.2]
    \coordinate (A) at (2,3);
    \coordinate (B) at (0,0);
    \coordinate (C) at (4,0);
    
    \draw[thick] (A) -- (B) -- (C) -- cycle;
    
    % Brūkšneliai kraštinėms (lygybės žymėjimas)
    \draw[thick, blue] (0.91, 1.56) -- (1.09, 1.44);
    \draw[thick, blue] (2.91, 1.44) -- (3.09, 1.56);
    
    % Kampai
    \pic [draw, fill=blue!10, fill opacity=0.35, "$\alpha$", angle radius=8mm,
           pic text options={text=black,text opacity=1,xshift=-1mm}] {angle = C--B--A};
    \pic [draw, fill=blue!10, fill opacity=0.35, "$\alpha$", angle radius=8mm,
           pic text options={text=black,text opacity=1,xshift=-1mm}] {angle = A--C--B};
    
    \foreach \p in {A,B,C} {\filldraw (\p) circle (1.2pt);}
    \node[above] at (A) {$A$};
    \node[below left] at (B) {$B$};
    \node[below right] at (C) {$C$};
\end{tikzpicture}
\end{center}
\textit{Išvada: Jei $AB = AC$, tai $\angle B = \angle C = \alpha$. Ir atvirkščiai!}

\subsection*{2. Greitkelis: Trikampio priekampis}
Mokiniai dažnai gaišta laiką skaičiuodami trečiąjį trikampio kampą ($180^\circ - \alpha - \beta$) ir tuomet atimdami jį iš $180^\circ$, kad gautų išorinį. \textbf{Tai neefektyvu!} Mokykitės matyti priekampius: trikampio priekampis (išorinis kampas) \textbf{iškart} lygus dviejų su juo negretimų vidaus kampų sumai.
\begin{center}
\begin{tikzpicture}[scale=1.1]
    \coordinate (A) at (1,2.5);
    \coordinate (B) at (0,0);
    \coordinate (C) at (3,0);
    \coordinate (D) at (5,0);
    
    \draw[thick] (A) -- (B) -- (D);
    \draw[thick] (A) -- (C);
    
    \pic [draw, fill=green!10, fill opacity=0.35, "$\beta$", angle radius=7mm,
        pic text options={text=black,text opacity=1}] {angle = B--A--C};
    \pic [draw, fill=blue!10, fill opacity=0.35, "$\alpha$", angle radius=7mm,
           pic text options={text=black,text opacity=1}] {angle=C--B--A};
        \pic [draw, thick, fill=red!20, fill opacity=0.35, "$\mathbf{\alpha+\beta}$", angle radius=8mm, angle eccentricity=1.5,
            pic text options={text=black,text opacity=1,xshift=-1mm}] {angle=D--C--A};
    
    \foreach \p in {A,B,C} {\filldraw (\p) circle (1.2pt);}
    \node[above] at (A) {$A$};
    \node[below left] at (B) {$B$};
    \node[below] at (C) {$C$};
\end{tikzpicture}
\end{center}

\subsection*{3. Euklido bėgiai: Lygiagrečių tiesių savybė}
Dvi lygiagrečias tieses perkirtus trečiąja, susidarantys vidaus priešiniai kampai yra lygūs.

Šiame pradžiamokslyje lygiagrečias tieses neoficialiai vadinsime \textbf{Euklido bėgiais}. Kai Euklido bėgiais pervažiuojame tiesę, gauname lygius vidaus priešinius kampus.
\begin{center}
\begin{tikzpicture}[scale=1.1]
    % Pagalbiniai taškai lygiagrečioms tiesėms (padeda išvengti klaidų žymint kampus)
    \coordinate (L1) at (0,2);
    \coordinate (R1) at (4,2);
    \coordinate (L2) at (0,0);
    \coordinate (R2) at (4,0);

    % Lygiagrečios tiesės
    \draw[thick, <->] (L1) -- (R1) node[right] {$a$};
    \draw[thick, <->] (L2) -- (R2) node[right] {$b \parallel a$};
    
    % Kertančioji
    \coordinate (P1) at (1,2.5);
    \coordinate (P2) at (3,-0.5);
    \draw[thick] (P1) -- (P2);
    
    \coordinate (Int1) at (1.33, 2);
    \coordinate (Int2) at (2.66, 0);
    
    % Z kampai (pataisyta priešlaikrodinė kryptis)
    \pic [draw, thick, fill=orange!20, fill opacity=0.35, "$\alpha$", angle radius=6mm, angle eccentricity=1.5,
        pic text options={text=black,text opacity=1}] {angle = L1--Int1--P2};
    \pic [draw, thick, fill=orange!20, fill opacity=0.35, "$\alpha$", angle radius=6mm, angle eccentricity=1.5,
        pic text options={text=black,text opacity=1}] {angle = R2--Int2--P1};
\end{tikzpicture}
\end{center}

% ==========================================
% PAVYZDŽIAI
% ==========================================
\section{Klasika: Kaip medžiojami kampai?}

Išnagrinėkime, kaip šios geometrinės savybės taikomos kartu: pirmiausia pasinaudosime lygiagrečių tiesių savybe, tuomet išvesime kampų lygybes, o galiausiai sudarysime ir išspręsime algebrinę lygtį.

\begin{example}[Trikampio kampų suma]
Įrodykte, jog bet kokio trikampio kampų suma lygi $180^\circ$.
\end{example}
\textbf{Sprendimas:}\\
Šią fundamentalią geometrijos tiesą lengvai įrodysime pasitelkę Euklido bėgiais.

\textbf{1 žingsnis.} Iš pradžių nusibraižome bet kokį trikampį $ABC$ ir jo vidinius kampus pažymime atitinkamai $\alpha$, $\beta$ bei $\gamma$.
\begin{center}
\begin{tikzpicture}[scale=1.1]
    \coordinate (A) at (1.2,2.8);
    \coordinate (B) at (0,0);
    \coordinate (C) at (4.5,0);

    \pic[draw, fill=green!20, fill opacity=0.35, "$\beta$", angle radius=6mm, pic text options={text=black,text opacity=1}] {angle=C--B--A};
    \pic[draw, fill=blue!20, fill opacity=0.35, "$\gamma$", angle radius=7mm, pic text options={text=black,text opacity=1,xshift=-1mm}] {angle=A--C--B};
    \pic[draw, fill=red!20, fill opacity=0.35, "$\alpha$", angle radius=6mm, pic text options={text=black,text opacity=1}] {angle=B--A--C};

    \draw[thick] (B) -- (A) -- (C) -- cycle;

    \foreach \p in {A,B,C} {\filldraw (\p) circle (1.2pt);}
    \node[above=2mm] at (A) {$A$};
    \node[below left] at (B) {$B$};
    \node[below right] at (C) {$C$};
\end{tikzpicture}
\end{center}

\textbf{2 žingsnis.} Euklido bėgiais pervažiuojame trikampį per viršūnę $A$ ir kraštinę $BC$: per viršūnę $A$ nubrėžiame pagalbinę tiesę $a$, lygiagrečią trikampio pagrindui $BC$. Kadangi $a \parallel BC$, o kraštinė $AB$ yra kirstinė, susidaro lygūs vidaus priešiniai kampai. Taigi kampo $\angle CBA$ porinis vidaus priešinis kampas prie viršūnės $A$ taip pat lygus $\beta$:
\begin{center}
\begin{tikzpicture}[scale=1.2]
    \coordinate (A) at (1.2,2.8);
    \coordinate (B) at (0,0);
    \coordinate (C) at (4.5,0);
    \coordinate (L) at (-1,2.8);
    \coordinate (R) at (5.5,2.8);

    \pic[draw, fill=green!20, fill opacity=0.35, "$\beta$", angle radius=6mm, pic text options={text=black,text opacity=1}] {angle=C--B--A};
    \pic[draw, fill=blue!20, fill opacity=0.35, "$\gamma$", angle radius=7mm, pic text options={text=black,text opacity=1,xshift=-1mm}] {angle=A--C--B};
    \pic[draw, fill=red!20, fill opacity=0.35, "$\alpha$", angle radius=6mm, pic text options={text=black,text opacity=1}] {angle=B--A--C};
    \pic[draw, fill=green!20, fill opacity=0.35, "$\beta$", angle radius=6mm, pic text options={text=black,text opacity=1}] {angle=L--A--B};

    \draw[thick] (B) -- (A) -- (C) -- cycle;
    \draw[thick, <->, blue] (L) -- (R) node[right, text=black] {$a \parallel BC$};
    \draw[thick, <->, blue] (-1,0) -- (5.5,0);

    \foreach \p in {A,B,C} {\filldraw (\p) circle (1.2pt);}
    \node[above=2mm] at (A) {$A$};
    \node[below left] at (B) {$B$};
    \node[below right] at (C) {$C$};
\end{tikzpicture}
\end{center}

\textbf{3 žingsnis.} Dešinėje pusėje kraštinė $AC$ taip pat veikia kaip kirstinė tarp tų pačių lygiagrečių tiesių. Todėl kampo $\angle BCA$ porinis vidaus priešinis kampas prie viršūnės $A$ lygus $\gamma$.
Matome, kad viršūnėje $A$ trys kampai $\beta$, $\alpha$ ir $\gamma$ susiglaudžia, sudarydami tiesę. Todėl šis ištiestinis kampas yra lygus $180^\circ$:
$$
\alpha + \beta + \gamma = 180^\circ
$$
\begin{center}
\begin{tikzpicture}[scale=1.3]
    \coordinate (A) at (1.2,2.8);
    \coordinate (B) at (0,0);
    \coordinate (C) at (4.5,0);
    \coordinate (L) at (-1,2.8);
    \coordinate (R) at (5.5,2.8);

    \pic[draw, fill=green!20, fill opacity=0.35, "$\beta$", angle radius=6mm, pic text options={text=black,text opacity=1}] {angle=C--B--A};
    \pic[draw, fill=blue!20, fill opacity=0.35, "$\gamma$", angle radius=7mm, pic text options={text=black,text opacity=1,xshift=-1mm}] {angle=A--C--B};
    \pic[draw, fill=red!20, fill opacity=0.35, "$\alpha$", angle radius=6mm, pic text options={text=black,text opacity=1}] {angle=B--A--C};
    \pic[draw, fill=green!20, fill opacity=0.35, "$\beta$", angle radius=6mm, pic text options={text=black,text opacity=1}] {angle=L--A--B};
    \pic[draw, fill=blue!20, fill opacity=0.35, "$\gamma$", angle radius=7mm, pic text options={text=black,text opacity=1,xshift=1mm,yshift=-0.5mm}] {angle=C--A--R};

    \draw[thick] (B) -- (A) -- (C) -- cycle;
    \draw[thick, <->, blue] (L) -- (R) node[right, text=black] {$a \parallel BC$};
    \draw[thick, <->, blue] (-1,0) -- (5.5,0);

    \foreach \p in {A,B,C} {\filldraw (\p) circle (1.2pt);}
    \node[above=2mm] at (A) {$A$};
    \node[below left] at (B) {$B$};
    \node[below right] at (C) {$C$};
\end{tikzpicture}
\end{center}
Įrodyta. \hfill $\square$


\newpage
\begin{example}[ „Bumerangas“]
Duotas neiškilusis keturkampis $ABCD$ (pavadinkime jį „bumerangu“), kurio trys vidiniai smailieji kampai yra $\angle A = \alpha$, $\angle B = \beta$ ir $\angle C = \gamma$. Įrodykite, kad jo išorinis (įgaubtasis) kampas $\angle ADC$ lygus šių trijų kampų sumai: $\alpha + \beta + \gamma$.
\end{example}
\textbf{Sprendimas:}\\
Šis uždavinys yra klasikinis pavyzdys, demonstruojantis trikampio priekampio savybės taikymą. Gautas rezultatas bus naudingas vėlesniame handoute apie apibrėžtinius keturkampius.

\textbf{1 žingsnis.} Tiksliai nubraižome uždavinio sąlygoje duotą neiškilųjį keturkampį $ABCD$ ir pusiau permatomomis spalvomis pažymime žinomus smailuosius kampus. Ieškomas įgaubtasis kampas $\angle ADC$ pažymėtas oranžine spalva ir klaustuku.
\begin{center}
\begin{tikzpicture}[scale=1.4]
    \coordinate (A) at (-2,-1);
    \coordinate (B) at (0,2);
    \coordinate (C) at (2,-1);
    \coordinate (D) at (0,0.2);

    \pic[draw, fill=blue!20, fill opacity=0.35, "$\alpha$", angle radius=6mm, angle eccentricity=1.4, pic text options={text=black, text opacity=1}] {angle=D--A--B};
    \pic[draw, fill=green!20, fill opacity=0.35, "$\beta$", angle radius=6mm, angle eccentricity=1.4, pic text options={text=black, text opacity=1}] {angle=A--B--C};
    \pic[draw, fill=purple!20, fill opacity=0.35, "$\gamma$", angle radius=6mm, angle eccentricity=1.4, pic text options={text=black, text opacity=1}] {angle=B--C--D};
    \pic[draw, fill=orange!20, fill opacity=0.35, "?", angle radius=5mm, angle eccentricity=1.6, pic text options={text=black, text opacity=1}] {angle=A--D--C};

    \draw[thick] (A) -- (B) -- (C) -- (D) -- cycle;

    \foreach \p in {A,B,C,D} {\filldraw (\p) circle (1.2pt);}
    \node[left] at (A) {$A$};
    \node[above] at (B) {$B$};
    \node[right] at (C) {$C$};
    \node[above left] at (D) {$D$};
\end{tikzpicture}
\end{center}

\textbf{2 žingsnis.} Pratęsiame atkarpą $AD$ tiesia linija, kol ji kerta priešingą kraštinę $BC$ taške $E$. „Bumerangas“ dabar padalintas į du trikampius: $\triangle ABE$ ir $\triangle CDE$.
\begin{center}
\begin{tikzpicture}[scale=1.4]
    \coordinate (A) at (-2,-1);
    \coordinate (B) at (0,2);
    \coordinate (C) at (2,-1);
    \coordinate (D) at (0,0.2);
    \coordinate (E) at (intersection of A--D and B--C);

    \pic[draw, fill=blue!20, fill opacity=0.35, "$\alpha$", angle radius=6mm, angle eccentricity=1.4, pic text options={text=black, text opacity=1}] {angle=D--A--B};
    \pic[draw, fill=green!20, fill opacity=0.35, "$\beta$", angle radius=6mm, angle eccentricity=1.4, pic text options={text=black, text opacity=1}] {angle=A--B--C};
    \pic[draw, fill=purple!20, fill opacity=0.35, "$\gamma$", angle radius=6mm, angle eccentricity=1.4, pic text options={text=black, text opacity=1}] {angle=B--C--D};
    \pic[draw, fill=orange!20, fill opacity=0.35, "?", angle radius=5mm, angle eccentricity=1.6, pic text options={text=black, text opacity=1}] {angle=A--D--C};

    \draw[thick] (A) -- (B) -- (C) -- (D) -- cycle;
    \draw[thick, dashed, blue] (D) -- (E);

    \foreach \p in {A,B,C,D,E} {\filldraw (\p) circle (1.2pt);}
    \node[left] at (A) {$A$};
    \node[above] at (B) {$B$};
    \node[right] at (C) {$C$};
    \node[above left] at (D) {$D$};
    \node[right] at (E) {$E$};
\end{tikzpicture}
\end{center}

\textbf{3 žingsnis.} Trikampyje $\triangle ABE$ kraštinė $BE$ yra pratęsta tolyn (iki taško $C$). Vadinasi, kampas $\angle AEC$ yra šio trikampio išorinis kampas (priekampis). Pagal priekampio savybę, jis lygus dviejų su juo negretimų vidaus kampų sumai:
$$
\angle AEC = \angle A + \angle B = \alpha + \beta
$$
Kadangi taškai $A, D, E$ guli vienoje tiesėje, tai $\angle DEC = \angle AEC$.
\begin{center}
\begin{tikzpicture}[scale=1.4]
    \coordinate (A) at (-2,-1);
    \coordinate (B) at (0,2);
    \coordinate (C) at (2,-1);
    \coordinate (D) at (0,0.2);
    \coordinate (E) at (intersection of A--D and B--C);

    \pic[draw, fill=blue!20, fill opacity=0.35, "$\alpha$", angle radius=6mm, angle eccentricity=1.4, pic text options={text=black, text opacity=1}] {angle=D--A--B};
    \pic[draw, fill=green!20, fill opacity=0.35, "$\beta$", angle radius=6mm, angle eccentricity=1.4, pic text options={text=black, text opacity=1}] {angle=A--B--C};
    \pic[draw, fill=purple!20, fill opacity=0.35, "$\gamma$", angle radius=6mm, angle eccentricity=1.4, pic text options={text=black, text opacity=1}] {angle=B--C--D};
    \pic[draw, fill=yellow!35, fill opacity=0.4, "$\alpha+\beta$", angle radius=5mm, angle eccentricity=1.8, pic text options={text=black, text opacity=1, xshift=1.5mm, yshift=-0.5mm}] {angle=A--E--C};
    \pic[draw, fill=orange!20, fill opacity=0.35, "?", angle radius=5mm, angle eccentricity=1.6, pic text options={text=black, text opacity=1}] {angle=A--D--C};

    \draw[thick] (A) -- (B) -- (C) -- (D) -- cycle;
    \draw[thick, dashed, blue] (D) -- (E);

    \foreach \p in {A,B,C,D,E} {\filldraw (\p) circle (1.2pt);}
    \node[left] at (A) {$A$};
    \node[above] at (B) {$B$};
    \node[right] at (C) {$C$};
    \node[above left] at (D) {$D$};
    \node[right] at (E) {$E$};
\end{tikzpicture}
\end{center}

\textbf{4 žingsnis.} Nagrinėkime mažąjį trikampį $\triangle CDE$. Jo atkarpa $ED$ yra pratęsta tolyn (iki taško $A$), todėl ieškomas įgaubtas kampas $\angle ADC$ yra ne kas kita, kaip trikampio $\triangle CDE$ išorinis kampas (priekampis) prie viršūnės $D$. Pritaikę priekampio savybę, gauname:
$$
\angle ADC = \angle DCE + \angle DEC = \gamma + (\alpha + \beta) = \alpha + \beta + \gamma
$$
\begin{center}
\begin{tikzpicture}[scale=1.4]
    \coordinate (A) at (-2,-1);
    \coordinate (B) at (0,2);
    \coordinate (C) at (2,-1);
    \coordinate (D) at (0,0.2);
    \coordinate (E) at (intersection of A--D and B--C);

    \pic[draw, fill=blue!20, fill opacity=0.35, "$\alpha$", angle radius=6mm, angle eccentricity=1.4, pic text options={text=black, text opacity=1}] {angle=D--A--B};
    \pic[draw, fill=green!20, fill opacity=0.35, "$\beta$", angle radius=6mm, angle eccentricity=1.4, pic text options={text=black, text opacity=1}] {angle=A--B--C};
    \pic[draw, fill=purple!20, fill opacity=0.35, "$\gamma$", angle radius=6mm, angle eccentricity=1.4, pic text options={text=black, text opacity=1}] {angle=B--C--D};
    \pic[draw, fill=yellow!35, fill opacity=0.4, "$\alpha+\beta$", angle radius=5mm, angle eccentricity=1.8, pic text options={text=black, text opacity=1, xshift=1.5mm, yshift=-0.5mm}] {angle=A--E--C};
    \pic[draw, fill=orange!25, fill opacity=0.45, "$\mathbf{\alpha+\beta+\gamma}$", angle radius=5mm, angle eccentricity=1.9, pic text options={text=black, text opacity=1}] {angle=A--D--C};

    \draw[thick] (A) -- (B) -- (C) -- (D) -- cycle;
    \draw[thick, dashed, blue] (D) -- (E);

    \foreach \p in {A,B,C,D,E} {\filldraw (\p) circle (1.2pt);}
    \node[left] at (A) {$A$};
    \node[above] at (B) {$B$};
    \node[right] at (C) {$C$};
    \node[above left] at (D) {$D$};
    \node[right] at (E) {$E$};
\end{tikzpicture}
\end{center}
Įrodyta. \hfill $\square$


\newpage
\begin{example}[Auksinis trikampis]
Duotas lygiašonis trikampis $ABC$, kurio kraštinės $AB = AC$. Kraštinėje $AC$ atidėtas taškas $D$ taip, kad $AD = BD = BC$. 
Įrodykite, kad $\angle A = 36^\circ$.
\end{example}

\textbf{Sprendimas:}\\
Spręsdami šį uždavinį pradėsime nuo vieno nežinomo kampo ir, nuosekliai naudodamiesi atkarpų lygybėmis, išreikšime visus kitus brėžinio kampus per tą patį nežinomąjį.

\textbf{1 žingsnis.} Nusibraižome trikampį $ABC$ ir kraštinėje $AC$ pažymime tašką $D$. Vienodais brūkšneliais atidedame sąlygoje nurodytas lygias atkarpas. Pirmajai lygių atkarpų grupei ($AD = BD = BC$) naudojame po vieną mėlyną brūkšnelį, o antrajai ($AB = AC$) – po du žalius brūkšnelius.
\begin{center}
\begin{tikzpicture}[scale=1.6]
    \coordinate (B) at (0,0);
    \coordinate (C) at (2,0);
    \coordinate (A) at (1,3.077);
    \coordinate (D) at ($(A)!0.618!(C)$);

    \draw[thick] (A) -- (B) -- (C) -- cycle;
    \draw[thick] (B) -- (D);

    % Vienodi brūkšneliai pirmai grupei
    \draw[thick, blue] ($(A)!0.5!(D) + (-0.095, -0.031)$) -- ($(A)!0.5!(D) + (0.095, 0.031)$);
    \draw[thick, blue] ($(B)!0.5!(D) + (-0.045, 0.089)$) -- ($(B)!0.5!(D) + (0.045, -0.089)$);
    \draw[thick, blue] (1, -0.1) -- (1, 0.1);

    % Du brūkšneliai antrai grupei
    \coordinate (mAB) at ($(A)!0.5!(B)$);
    \draw[thick, green!60!black] ($(mAB) + (-0.08, 0.04)$) -- ($(mAB) + (0.08, -0.04)$);
    \draw[thick, green!60!black] ($(mAB) + (-0.04, 0.09)$) -- ($(mAB) + (0.12, 0.01)$);
    \coordinate (mAC) at ($(A)!0.5!(C)$);
    \draw[thick, green!60!black] ($(mAC) + (-0.08, -0.04)$) -- ($(mAC) + (0.08, 0.04)$);
    \draw[thick, green!60!black] ($(mAC) + (-0.04, -0.09)$) -- ($(mAC) + (0.12, -0.01)$);

    \foreach \p in {A,B,C,D} {\filldraw (\p) circle (1.2pt);}
    \node[above] at (A) {$A$};
    \node[below left] at (B) {$B$};
    \node[below right] at (C) {$C$};
    \node[right] at (D) {$D$};
\end{tikzpicture}
\end{center}

\textbf{2 žingsnis.} Pažymime viršūnės kampą $\angle A = \alpha$. Kadangi atkarpos $AD = BD$, trikampis $\triangle ADB$ yra lygiašonis. Todėl jo kampai prie pagrindo $AB$ privalo būti lygūs:
$$
\angle ABD = \angle A = \alpha
$$
\begin{center}
\begin{tikzpicture}[scale=1.6]
    \coordinate (B) at (0,0); \coordinate (C) at (2,0); \coordinate (A) at (1,3.077); \coordinate (D) at ($(A)!0.618!(C)$);

    \pic[draw, fill=blue!20, fill opacity=0.35, "$\alpha$", angle radius=6mm, angle eccentricity=1.4, pic text options={text=black, text opacity=1}] {angle=B--A--C};
    \pic[draw, fill=blue!20, fill opacity=0.35, "$\alpha$", angle radius=6mm, angle eccentricity=1.4, pic text options={text=black, text opacity=1, xshift=2mm, yshift=1mm}] {angle=D--B--A};

    \draw[thick] (A) -- (B) -- (C) -- cycle; \draw[thick] (B) -- (D);

    \draw[thick, blue] ($(A)!0.5!(D) + (-0.095, -0.031)$) -- ($(A)!0.5!(D) + (0.095, 0.031)$);
    \draw[thick, blue] ($(B)!0.5!(D) + (-0.045, 0.089)$) -- ($(B)!0.5!(D) + (0.045, -0.089)$);
    \draw[thick, blue] (1, -0.1) -- (1, 0.1);
    \coordinate (mAB) at ($(A)!0.5!(B)$);
    \draw[thick, green!60!black] ($(mAB) + (-0.08, 0.04)$) -- ($(mAB) + (0.08, -0.04)$);
    \draw[thick, green!60!black] ($(mAB) + (-0.04, 0.09)$) -- ($(mAB) + (0.12, 0.01)$);
    \coordinate (mAC) at ($(A)!0.5!(C)$);
    \draw[thick, green!60!black] ($(mAC) + (-0.08, -0.04)$) -- ($(mAC) + (0.08, 0.04)$);
    \draw[thick, green!60!black] ($(mAC) + (-0.04, -0.09)$) -- ($(mAC) + (0.12, -0.01)$);

    \foreach \p in {A,B,C,D} {\filldraw (\p) circle (1.2pt);}
    \node[above] at (A) {$A$}; \node[below left] at (B) {$B$}; \node[below right] at (C) {$C$}; \node[right] at (D) {$D$};
\end{tikzpicture}
\end{center}

\textbf{3 žingsnis.} Kampas $\angle BDC$ yra trikampio $\triangle ADB$ išorinis kampas (priekampis) prie viršūnės $D$. Pagal priekampio savybę jis lygus dviejų su juo negretimų vidaus kampų sumai:
$$
\angle BDC = \angle A + \angle ABD = \alpha + \alpha = 2\alpha
$$
\begin{center}
\begin{tikzpicture}[scale=1.6]
    \coordinate (B) at (0,0); \coordinate (C) at (2,0); \coordinate (A) at (1,3.077); \coordinate (D) at ($(A)!0.618!(C)$);

    \pic[draw, fill=blue!20, fill opacity=0.35, "$\alpha$", angle radius=6mm, angle eccentricity=1.4, pic text options={text=black, text opacity=1}] {angle=B--A--C};
    \pic[draw, fill=blue!20, fill opacity=0.35, "$\alpha$", angle radius=6mm, angle eccentricity=1.4, pic text options={text=black, text opacity=1, yshift=2mm}] {angle=D--B--A};
    \pic[draw, fill=red!20, fill opacity=0.35, "$2\alpha$", angle radius=5mm, angle eccentricity=1.5, pic text options={text=black, text opacity=1}] {angle=B--D--C};

    \draw[thick] (A) -- (B) -- (C) -- cycle; \draw[thick] (B) -- (D);

    \draw[thick, blue] ($(A)!0.5!(D) + (-0.095, -0.031)$) -- ($(A)!0.5!(D) + (0.095, 0.031)$);
    \draw[thick, blue] ($(B)!0.5!(D) + (-0.045, 0.089)$) -- ($(B)!0.5!(D) + (0.045, -0.089)$);
    \draw[thick, blue] (1, -0.1) -- (1, 0.1);
    \coordinate (mAB) at ($(A)!0.5!(B)$);
    \draw[thick, green!60!black] ($(mAB) + (-0.08, 0.04)$) -- ($(mAB) + (0.08, -0.04)$);
    \draw[thick, green!60!black] ($(mAB) + (-0.04, 0.09)$) -- ($(mAB) + (0.12, 0.01)$);
    \coordinate (mAC) at ($(A)!0.5!(C)$);
    \draw[thick, green!60!black] ($(mAC) + (-0.08, -0.04)$) -- ($(mAC) + (0.08, 0.04)$);
    \draw[thick, green!60!black] ($(mAC) + (-0.04, -0.09)$) -- ($(mAC) + (0.12, -0.01)$);

    \foreach \p in {A,B,C,D} {\filldraw (\p) circle (1.2pt);}
    \node[above] at (A) {$A$}; \node[below left] at (B) {$B$}; \node[below right] at (C) {$C$}; \node[right] at (D) {$D$};
\end{tikzpicture}
\end{center}

\textbf{4 žingsnis.} Pagal sąlygą turime ir $BD = BC$, vadinasi, mažasis trikampis $\triangle BDC$ taip pat yra lygiašonis. Jo kampai prie pagrindo $DC$ yra lygūs, todėl:
$$
\angle C = \angle BDC = 2\alpha
$$
\begin{center}
\begin{tikzpicture}[scale=1.6]
    \coordinate (B) at (0,0); \coordinate (C) at (2,0); \coordinate (A) at (1,3.077); \coordinate (D) at ($(A)!0.618!(C)$);

    \pic[draw, fill=blue!20, fill opacity=0.35, "$\alpha$", angle radius=6mm, angle eccentricity=1.4, pic text options={text=black, text opacity=1}] {angle=B--A--C};
    \pic[draw, fill=blue!20, fill opacity=0.35, "$\alpha$", angle radius=6mm, angle eccentricity=1.4, pic text options={text=black, text opacity=1}] {angle=D--B--A};
    \pic[draw, fill=red!20, fill opacity=0.35, "$2\alpha$", angle radius=5mm, angle eccentricity=1.5, pic text options={text=black, text opacity=1}] {angle=B--D--C};
    \pic[draw, fill=red!20, fill opacity=0.35, "$2\alpha$", angle radius=5mm, angle eccentricity=1.5, pic text options={text=black, text opacity=1}] {angle=A--C--B};

    \draw[thick] (A) -- (B) -- (C) -- cycle; \draw[thick] (B) -- (D);

    \draw[thick, blue] ($(A)!0.5!(D) + (-0.095, -0.031)$) -- ($(A)!0.5!(D) + (0.095, 0.031)$);
    \draw[thick, blue] ($(B)!0.5!(D) + (-0.045, 0.089)$) -- ($(B)!0.5!(D) + (0.045, -0.089)$);
    \draw[thick, blue] (1, -0.1) -- (1, 0.1);
    \coordinate (mAB) at ($(A)!0.5!(B)$);
    \draw[thick, green!60!black] ($(mAB) + (-0.08, 0.04)$) -- ($(mAB) + (0.08, -0.04)$);
    \draw[thick, green!60!black] ($(mAB) + (-0.04, 0.09)$) -- ($(mAB) + (0.12, 0.01)$);
    \coordinate (mAC) at ($(A)!0.5!(C)$);
    \draw[thick, green!60!black] ($(mAC) + (-0.08, -0.04)$) -- ($(mAC) + (0.08, 0.04)$);
    \draw[thick, green!60!black] ($(mAC) + (-0.04, -0.09)$) -- ($(mAC) + (0.12, -0.01)$);

    \foreach \p in {A,B,C,D} {\filldraw (\p) circle (1.2pt);}
    \node[above] at (A) {$A$}; \node[below left] at (B) {$B$}; \node[below right] at (C) {$C$}; \node[right] at (D) {$D$};
\end{tikzpicture}
\end{center}

\textbf{5 žingsnis.} Grįžtame prie didžiojo trikampio $ABC$. Dvigubi žali brūkšneliai rodo, kad $AB = AC$, vadinasi, ir jis yra lygiašonis. Todėl visas kampas prie viršūnės $B$ privalo būti lygus kampui prie viršūnės $C$:
$$
\angle ABC = \angle C = 2\alpha
$$
\begin{center}
\begin{tikzpicture}[scale=1.6]
    \coordinate (B) at (0,0); \coordinate (C) at (2,0); \coordinate (A) at (1,3.077); \coordinate (D) at ($(A)!0.618!(C)$);

    \pic[draw, fill=blue!20, fill opacity=0.35, "$\alpha$", angle radius=6mm, angle eccentricity=1.4, pic text options={text=black, text opacity=1}] {angle=B--A--C};
    \pic[draw, fill=blue!20, fill opacity=0.35, "$\alpha$", angle radius=6mm, angle eccentricity=1.4, pic text options={text=black, text opacity=1, yshift=2mm}] {angle=D--B--A};
    \pic[draw, fill=red!20, fill opacity=0.35, "$2\alpha$", angle radius=5mm, angle eccentricity=1.5, pic text options={text=black, text opacity=1}] {angle=B--D--C};
    \pic[draw, fill=red!20, fill opacity=0.35, "$2\alpha$", angle radius=5mm, angle eccentricity=1.5, pic text options={text=black, text opacity=1}] {angle=A--C--B};
    
    % Visas didysis kampas ABC
    \pic[draw, fill=red!20, fill opacity=0.35, "$2\alpha$", angle radius=8mm, angle eccentricity=1.15, pic text options={text=black, text opacity=1, xshift=2.5mm, yshift=-1.5mm}] {angle=C--B--A};

    \draw[thick] (A) -- (B) -- (C) -- cycle; \draw[thick] (B) -- (D);

    \draw[thick, blue] ($(A)!0.5!(D) + (-0.095, -0.031)$) -- ($(A)!0.5!(D) + (0.095, 0.031)$);
    \draw[thick, blue] ($(B)!0.5!(D) + (-0.045, 0.089)$) -- ($(B)!0.5!(D) + (0.045, -0.089)$);
    \draw[thick, blue] (1, -0.1) -- (1, 0.1);
    \coordinate (mAB) at ($(A)!0.5!(B)$);
    \draw[thick, green!60!black] ($(mAB) + (-0.08, 0.04)$) -- ($(mAB) + (0.08, -0.04)$);
    \draw[thick, green!60!black] ($(mAB) + (-0.04, 0.09)$) -- ($(mAB) + (0.12, 0.01)$);
    \coordinate (mAC) at ($(A)!0.5!(C)$);
    \draw[thick, green!60!black] ($(mAC) + (-0.08, -0.04)$) -- ($(mAC) + (0.08, 0.04)$);
    \draw[thick, green!60!black] ($(mAC) + (-0.04, -0.09)$) -- ($(mAC) + (0.12, -0.01)$);

    \foreach \p in {A,B,C,D} {\filldraw (\p) circle (1.2pt);}
    \node[above] at (A) {$A$}; \node[below left] at (B) {$B$}; \node[below right] at (C) {$C$}; \node[right] at (D) {$D$};
\end{tikzpicture}
\end{center}

\textbf{6 žingsnis.} Dabar visi trys didžiojo trikampio $\triangle ABC$ kampai išreikšti dydžiu $\alpha$. Kadangi bet kurio trikampio vidaus kampų suma visada lygi $180^\circ$, gauname:
$$
\angle A + \angle ABC + \angle C = 180^\circ
$$
$$
\alpha + 2\alpha + 2\alpha = 180^\circ \implies 5\alpha = 180^\circ \implies \alpha = 36^\circ
$$
Įrodyta. \hfill $\square$


\newpage
\begin{example}[Kampas tarp pusiaukampinių]
Trikampio $ABC$ kampas $\angle A = \alpha$. Iš viršūnių $B$ ir $C$ nubrėžtos pusiaukampinės susikerta taške $I$. Įrodykite, kad kampas tarp jų $\angle BIC = 90^\circ + \frac{\alpha}{2}$.
\end{example}
\textbf{Sprendimas:}\\
Šis pavyzdys tobulai iliustruoja atvirkštinį loginį mąstymą: pradedame nuo to, ko mūsų prašo, ir tuomet ieškome trūkstamų duomenų.

\textbf{1 žingsnis.} Nusibraižome trikampį $ABC$ ir jo viduje taške $I$ susikertančias pusiaukampines $BI$ bei $CI$. Viršūnės kampą pažymime mėlyna spalva ir $\alpha$. Nors atkarpas žinome esant pusiaukampinėmis, pagal taisykles išankstinių lygybių brėžinyje nežymime. Ieškomą kampą $\angle BIC$ pažymime oranžine spalva su klaustuku.
\begin{center}
\begin{tikzpicture}[scale=1.2]
    \coordinate (B) at (0,0);
    \coordinate (C) at (5,0);
    \coordinate (A) at (1.5,3.5);
    \coordinate (I) at (2.2, 1.2);

    \pic[draw, fill=blue!20, fill opacity=0.35, "$\alpha$", angle radius=6mm, pic text options={text=black,text opacity=1}] {angle = B--A--C};
    \pic[draw, fill=orange!20, fill opacity=0.4, "?", angle radius=5mm, angle eccentricity=1.5, pic text options={text=black,text opacity=1}] {angle = B--I--C};

    \draw[thick] (A) -- (B) -- (C) -- cycle;
    \draw[thick, blue] (B) -- (I) -- (C);

    \foreach \p in {A,B,C,I} {\filldraw (\p) circle (1.2pt);}
    \node[above] at (A) {$A$};
    \node[below left] at (B) {$B$};
    \node[below right] at (C) {$C$};
    \node[above, yshift=1mm] at (I) {$I$};
\end{tikzpicture}
\end{center}

\textbf{2 žingsnis.} Kadangi $BI$ ir $CI$ yra pusiaukampinės, jos atitinkamus kampus padalija per pusę. Pažymėkime susidariusias puses vienodomis spalvomis ir raidėmis: kampui $B$ priskirkime $\angle IBC = \angle IBA = \beta$ (geltona), o kampui $C$ – $\angle ICB = \angle ICA = \gamma$ (žalia).
\begin{center}
\begin{tikzpicture}[scale=1.2]
    \coordinate (B) at (0,0);
    \coordinate (C) at (5,0);
    \coordinate (A) at (1.5,3.5);
    \coordinate (I) at (2.2, 1.2);

    \pic[draw, fill=blue!20, fill opacity=0.35, "$\alpha$", angle radius=6mm, pic text options={text=black,text opacity=1}] {angle = B--A--C};
    \pic[draw, fill=yellow!25, fill opacity=0.35, "$\beta$", angle radius=6mm, angle eccentricity=1.4, pic text options={text=black,text opacity=1}] {angle = C--B--I};
    \pic[draw, fill=yellow!25, fill opacity=0.35, "$\beta$", angle radius=8mm, angle eccentricity=1.25, pic text options={text=black,text opacity=1}] {angle = I--B--A};
    \pic[draw, fill=green!20, fill opacity=0.35, "$\gamma$", angle radius=7mm, angle eccentricity=1.35, pic text options={text=black,text opacity=1}] {angle = A--C--I};
    \pic[draw, fill=green!20, fill opacity=0.35, "$\gamma$", angle radius=9mm, angle eccentricity=1.2, pic text options={text=black,text opacity=1}] {angle = I--C--B};
    \pic[draw, fill=orange!20, fill opacity=0.4, "?", angle radius=5mm, angle eccentricity=1.5, pic text options={text=black,text opacity=1}] {angle = B--I--C};

    \draw[thick] (A) -- (B) -- (C) -- cycle;
    \draw[thick, blue] (B) -- (I) -- (C);

    \foreach \p in {A,B,C,I} {\filldraw (\p) circle (1.2pt);}
    \node[above] at (A) {$A$};
    \node[below left] at (B) {$B$};
    \node[below right] at (C) {$C$};
    \node[above, yshift=1mm] at (I) {$I$};
\end{tikzpicture}
\end{center}

\textbf{3 žingsnis.} Nagrinėkime mažąjį trikampį $\triangle BIC$. Jo kampų suma lygi $180^\circ$. Iš šios taisyklės galime išreikšti mūsų ieškomą kampą $\angle BIC$:
$$
\angle BIC + \beta + \gamma = 180^\circ \implies \angle BIC = 180^\circ - (\beta + \gamma)
$$
Matome, kad norint sužinoti ieškomą kampą, pakanka sužinoti sumą $(\beta + \gamma)$.
\begin{center}
\begin{tikzpicture}[scale=1.2]
    \coordinate (B) at (0,0);
    \coordinate (C) at (5,0);
    \coordinate (A) at (1.5,3.5);
    \coordinate (I) at (2.2, 1.2);

    \pic[draw, fill=blue!20, fill opacity=0.35, "$\alpha$", angle radius=6mm, pic text options={text=black,text opacity=1}] {angle = B--A--C};
    \pic[draw, fill=yellow!25, fill opacity=0.35, "$\beta$", angle radius=6mm, angle eccentricity=1.4, pic text options={text=black,text opacity=1}] {angle = C--B--I};
    \pic[draw, fill=yellow!25, fill opacity=0.35, "$\beta$", angle radius=8mm, angle eccentricity=1.25, pic text options={text=black,text opacity=1}] {angle = I--B--A};
    \pic[draw, fill=green!20, fill opacity=0.35, "$\gamma$", angle radius=7mm, angle eccentricity=1.35, pic text options={text=black,text opacity=1}] {angle = A--C--I};
    \pic[draw, fill=green!20, fill opacity=0.35, "$\gamma$", angle radius=9mm, angle eccentricity=1.2, pic text options={text=black,text opacity=1}] {angle = I--C--B};
    \pic[draw, fill=orange!20, fill opacity=0.4, "?", angle radius=5mm, angle eccentricity=1.5, pic text options={text=black,text opacity=1}] {angle = B--I--C};

    \draw[thick] (A) -- (B) -- (C) -- cycle;
    \draw[thick, blue] (B) -- (I) -- (C);

    \foreach \p in {A,B,C,I} {\filldraw (\p) circle (1.2pt);}
    \node[above] at (A) {$A$};
    \node[below left] at (B) {$B$};
    \node[below right] at (C) {$C$};
    \node[above, yshift=1mm] at (I) {$I$};
\end{tikzpicture}
\end{center}

\textbf{4 žingsnis.} Dabar pažiūrėkime į visą didįjį trikampį $\triangle ABC$. Jo kampų suma taip pat yra $180^\circ$. Išreiškę visus kampus naudodamiesi pažymėtomis raidėmis turime:
$$
\alpha + 2\beta + 2\gamma = 180^\circ
$$
Iškelkime dvejetą prieš skliaustus ir iš lygties išreikškime mums reikiamą sumą $(\beta + \gamma)$:
$$
2(\beta + \gamma) = 180^\circ - \alpha \implies \beta + \gamma = 90^\circ - \frac{\alpha}{2}
$$
\begin{center}
\begin{tikzpicture}[scale=1.2]
    \coordinate (B) at (0,0);
    \coordinate (C) at (5,0);
    \coordinate (A) at (1.5,3.5);
    \coordinate (I) at (2.2, 1.2);

    \pic[draw, fill=blue!20, fill opacity=0.35, "$\alpha$", angle radius=6mm, pic text options={text=black,text opacity=1}] {angle = B--A--C};
    \pic[draw, fill=yellow!25, fill opacity=0.35, "$\beta$", angle radius=6mm, angle eccentricity=1.4, pic text options={text=black,text opacity=1}] {angle = C--B--I};
    \pic[draw, fill=yellow!25, fill opacity=0.35, "$\beta$", angle radius=8mm, angle eccentricity=1.25, pic text options={text=black,text opacity=1}] {angle = I--B--A};
    \pic[draw, fill=green!20, fill opacity=0.35, "$\gamma$", angle radius=7mm, angle eccentricity=1.35, pic text options={text=black,text opacity=1}] {angle = A--C--I};
    \pic[draw, fill=green!20, fill opacity=0.35, "$\gamma$", angle radius=9mm, angle eccentricity=1.2, pic text options={text=black,text opacity=1}] {angle = I--C--B};
    \pic[draw, fill=orange!20, fill opacity=0.4, "?", angle radius=5mm, angle eccentricity=1.5, pic text options={text=black,text opacity=1}] {angle = B--I--C};

    \draw[thick] (A) -- (B) -- (C) -- cycle;
    \draw[thick, blue] (B) -- (I) -- (C);

    \foreach \p in {A,B,C,I} {\filldraw (\p) circle (1.2pt);}
    \node[above] at (A) {$A$};
    \node[below left] at (B) {$B$};
    \node[below right] at (C) {$C$};
    \node[above, yshift=1mm] at (I) {$I$};
\end{tikzpicture}
\end{center}

\textbf{5 žingsnis.} Grįžtame prie 3-ajame žingsnyje sudarytos lygties ieškomam kampui $\angle BIC$ ir vietoje bloko $(\beta + \gamma)$ įstatome ką tik atrastą jo išraišką:
$$
\angle BIC = 180^\circ - \left( 90^\circ - \frac{\alpha}{2} \right) = 180^\circ - 90^\circ + \frac{\alpha}{2} = 90^\circ + \frac{\alpha}{2}
$$
\begin{center}
\begin{tikzpicture}[scale=1.2]
    \coordinate (B) at (0,0);
    \coordinate (C) at (5,0);
    \coordinate (A) at (1.5,3.5);
    \coordinate (I) at (2.2, 1.2);

    \pic[draw, fill=blue!20, fill opacity=0.35, "$\alpha$", angle radius=6mm, pic text options={text=black,text opacity=1}] {angle = B--A--C};
    \pic[draw, fill=yellow!25, fill opacity=0.35, "$\beta$", angle radius=6mm, angle eccentricity=1.4, pic text options={text=black,text opacity=1}] {angle = C--B--I};
    \pic[draw, fill=yellow!25, fill opacity=0.35, "$\beta$", angle radius=8mm, angle eccentricity=1.25, pic text options={text=black,text opacity=1}] {angle = I--B--A};
    \pic[draw, fill=green!20, fill opacity=0.35, "$\gamma$", angle radius=7mm, angle eccentricity=1.35, pic text options={text=black,text opacity=1}] {angle = A--C--I};
    \pic[draw, fill=green!20, fill opacity=0.35, "$\gamma$", angle radius=9mm, angle eccentricity=1.2, pic text options={text=black,text opacity=1}] {angle = I--C--B};
    \pic[draw, fill=orange!25, fill opacity=0.45, "$90^\circ + \frac{\alpha}{2}$", angle radius=5mm, angle eccentricity=2.2, pic text options={text=black,text opacity=1}] {angle = B--I--C};

    \draw[thick] (A) -- (B) -- (C) -- cycle;
    \draw[thick, blue] (B) -- (I) -- (C);

    \foreach \p in {A,B,C,I} {\filldraw (\p) circle (1.2pt);}
    \node[above] at (A) {$A$};
    \node[below left] at (B) {$B$};
    \node[below right] at (C) {$C$};
    \node[above, yshift=1mm] at (I) {$I$};
\end{tikzpicture}
\end{center}
Įrodyta. \hfill $\square$

\vspace{0.5cm}

% ==========================================
% UŽDAVINIAI
% ==========================================
\newpage
\section{Uždaviniai savarankiškam sprendimui (12 iššūkių)}

\begin{enumerate}
    \renewcommand{\labelenumi}{\textbf{\theenumi.}}
    
    \item Lygiašonio trikampio $ABC$ ($AB=BC$) išorinis kampas prie viršūnės $A$ lygus $140^\circ$. Apskaičiuokite trikampio $ABC$ kampus.
    \item Trikampio $ABC$ $\angle A = 40^\circ$, o $\angle B = 80^\circ$. Iš viršūnės $C$ nubrėžta pusiaukampinė $CD$. Apskaičiuokite $\angle BDC$.
    \item Trikampyje $ABC$ nubrėžta aukštinė $AH$ ir pusiaukampinė $AD$. Žinoma, kad $\angle B = 70^\circ$, o $\angle C = 40^\circ$. Apskaičiuokite kampą tarp aukštinės ir pusiaukampinės ($\angle HAD$).
    
    \item Tarp dviejų lygiagrečių tiesių $a$ ir $b$ nubrėžta laužtė iš dviejų atkarpų. Pažymėję prie tiesių esančius smailiuosius kampus $\alpha$ ir $\beta$, įrodykite, kad kampas laužtės „smaigalyje“ lygus $\alpha + \beta$.
    \begin{center}
    \begin{tikzpicture}[scale=1.2]
        \coordinate (A1) at (0,2); \coordinate (A2) at (5,2);
        \coordinate (B1) at (0,0); \coordinate (B2) at (5,0);
        
        \draw[thick, <->] (A1) -- (A2) node[right] {$a$};
        \draw[thick, <->] (B1) -- (B2) node[right] {$b \parallel a$};
        
        \coordinate (P) at (1.5,2);
        \coordinate (Q) at (1.712,0.8);
        \coordinate (R) at (1.041,0);
        
        \draw[thick] (P) -- (Q) -- (R);
        
        \pic[draw, fill=blue!10, "$\alpha$", angle radius=5mm, angle eccentricity=1.1,
            pic text options={text=black,text opacity=1,xshift=-2mm, yshift=-1mm}] {angle=A1--P--Q};
        \pic[draw, fill=green!10, "$\beta$", angle radius=5mm, angle eccentricity=1.1,
            pic text options={text=black,text opacity=1,xshift=2mm, yshift=1mm}] {angle=B2--R--Q};
        \pic[draw, fill=orange!20, fill opacity=0.35, "?", angle radius=5mm, angle eccentricity=1.3,
            pic text options={text=black,text opacity=1}] {angle=P--Q--R};
    \end{tikzpicture}
    \end{center}
    
    \item Per lygiašonio trikampio $ABC$ ($AB=AC$) viršūnę $A$ nubrėžta tiesė, lygiagreti pagrindui $BC$. Įrodykite, kad trikampio šoninės kraštinės (arba jų tęsiniai) su šia tiese sudaro lygius kampus.
    \item Trikampyje $ABC$ nubrėžta pusiaukampinė $AK$. Per tašką $K$ nubrėžta tiesė, lygiagreti kraštinei $AB$, kerta kraštinę $AC$ taške $M$. Įrodykite, kad trikampis $AMK$ lygiašonis.
    
    \item Penkiakampė žvaigždė: Nubraižykite bet kokią penkiakampę žvaigždę (sudarytą iš 5 atkarpų). Įrodykite, kad visų penkių jos „smaigalių“ kampų suma lygi $180^\circ$.
    \begin{center}
    \begin{tikzpicture}[scale=1.1]
        \coordinate (P1) at (90:2);
        \coordinate (P2) at (18:2);
        \coordinate (P3) at (-54:2);
        \coordinate (P4) at (-126:2);
        \coordinate (P5) at (162:2);
        
        \draw[thick] (P1) -- (P3) -- (P5) -- (P2) -- (P4) -- cycle;
        
        \pic[draw, fill=red!20, "1", angle radius=5mm, pic text options={scale=0.8}] {angle=P4--P1--P3};
        \pic[draw, fill=green!20, "2", angle radius=5mm, pic text options={scale=0.8}] {angle=P5--P2--P4};
        \pic[draw, fill=blue!20, "3", angle radius=5mm, pic text options={scale=0.8}] {angle=P1--P3--P5};
        \pic[draw, fill=orange!20, "4", angle radius=5mm, pic text options={scale=0.8}] {angle=P2--P4--P1};
        \pic[draw, fill=purple!20, "5", angle radius=5mm, pic text options={scale=0.8}] {angle=P3--P5--P2};
    \end{tikzpicture}
    \end{center}
    
    \item Stačiojo trikampio $ABC$ ($\angle C = 90^\circ$) įžambinėje $AB$ atidėti taškai $M$ ir $N$ taip, kad $AM = AC$ ir $BN = BC$. Apskaičiuokite $\angle MCN$.
    \item Trikampio $ABC$ išorinių kampų prie viršūnių $B$ ir $C$ pusiaukampinės susikerta taške $J$. Remdamiesi 2-ojo pavyzdžio („Bumerango“ konstrukcijos) idėjomis, įrodykite, kad $\angle BJC = 90^\circ - \frac{\angle A}{2}$.
    \item Kvadrato $ABCD$ viduje nubrėžtas lygiakraštis trikampis $ABE$. Apskaičiuokite $\angle CDE$.
    
    \item Stačiajame trikampyje $ABC$ ($\angle C = 90^\circ$) nubrėžta pusiaukraštinė $CM$. Įrodykite, kad jos ilgis lygus pusei įžambinės ($CM = \frac{AB}{2}$).
    \item Trikampyje $ABC$ $\angle B = 20^\circ$, o $\angle C = 40^\circ$. Kraštinėje $BC$ atidėtas taškas $D$ taip, kad $AD = DC$. Apskaičiuokite $\angle BAD$.
\end{enumerate}

\end{document}